% GENERATED FILE. Edit canonical lessons, then run npm run generate:books.
% canonical-source-hash: fnv1a64:6f8df0c03d11807f
% canonical-lessons: ES-C295-sendero, ES-C295-puente, ES-C295-orilla, ES-C295-isla, ES-C295-niebla

\chapter{The Way Out --- Path, Bridge, Shore, Island, Mist}
\label{ch:el-sendero}
\hlchaptermodality{295}

\begin{chapteropening}
\textbf{By the end of this chapter:} \emph{I can name the path, the bridge, the shore, the island and the mist in Spanish.}

Complete ES-C295-niebla: name all five of the way out, from the path to the mist behind you.
\end{chapteropening}

\section[el sendero]{el sendero — a half-sized road}

\label{lesson:ES-C295-sendero}



\begin{quote}
Before a new run: what was \emph{scopa}? (A twig, and a bundle of them.) And what was \emph{mantum}? (A cloak.)
\end{quote}

\subsection*{What to know first}
\begin{itemize}
\raggedright
  \item la salida — where one of these begins.
  \item el bosque — what it usually runs through.
\end{itemize}

\begin{sounds}
\begin{itemize}
\raggedright
  \item \emph{sen-DE-ro}, stress in the middle. The \emph{d} between vowels is soft, and the \emph{r} is a single tap. $\to$ reference
\end{itemize}
\end{sounds}

\begin{cousinweb}
\textbf{el sendero} — \textquotedblleft{}the path,\textquotedblright{} \textquotedblleft{}the trail.\textquotedblright{}

Behind it is Latin \textbf{semita}, a footpath, and \emph{semita} looks to have been built from \textbf{semi-}, half, and a piece meaning going. A footpath was a half-road: narrower than the paved thing the legions marched on, and going somewhere less official.

That \emph{semi-} is the one you already know from \textbf{semicircle} and \textbf{semifinal}. It is doing exactly the same job here as it does there — halving something — and what it halves in \emph{sendero} is the road itself.

Spanish also has \emph{camino} for the larger sort of road. So the language kept the Roman distinction: one word for the way everybody takes, another for the way through the trees.
\end{cousinweb}

\begin{grammarlens}[title={What you've built}]
One word, and the start of a way out.

You have had \emph{la salida} since the departure words. A \emph{salida} is where you leave; a \emph{sendero} is what you leave along. The two fit together, and this run is going to keep fitting things onto the end of them.
\end{grammarlens}

\subsection*{Your turn}
Say these aloud:

\begin{itemize}
\raggedright
  \item \emph{el sendero}
  \item \emph{el sendero del bosque}
  \item \emph{la salida}, then \emph{el sendero}
\end{itemize}

\begin{tcolorbox}[breakable,colback=teal!4,colframe=teal!35!black,title={Before you move on}]
What was \emph{semita}? (A footpath.) What does the \emph{semi-} in it mean? (Half.) And what is \emph{el sendero}? (The path.)
\end{tcolorbox}
\section[el puente]{el puente — why the pope is a bridge-builder}

\label{lesson:ES-C295-puente}



\begin{quote}
Two back: what was \emph{scopa}? (A twig.) And what does the \emph{semi-} in \emph{sendero} mean? (Half.)
\end{quote}

\subsection*{What to know first}
\begin{itemize}
\raggedright
  \item el sendero — what this carries across.
  \item el lago — the sort of thing it goes over.
\end{itemize}

\begin{sounds}
\begin{itemize}
\raggedright
  \item \emph{PUEN-te}, stress on the first part. The \emph{ue} is one slide, the same as in \emph{pueblo} and \emph{cuello}. $\to$ reference
\end{itemize}
\end{sounds}

\begin{cousinweb}
\textbf{el puente} — \textquotedblleft{}the bridge.\textquotedblright{}

Latin \textbf{pons}, and in its other form \textbf{pontem}. There is the \emph{o}-to-\emph{ue} break for the third time — \emph{pueblo}, \emph{cuello}, and now \emph{puente}. By this point you should be able to predict it.

The English relatives are worth the detour. A \textbf{pontoon} is a floating bridge. To \textbf{pontificate} is to speak as a \textbf{pontiff} does — and \emph{pontifex}, the Roman priestly title the popes inherited, is \emph{pons} plus \emph{facere}: a bridge-maker. Nobody is certain whether the Roman priests literally maintained the bridge over the Tiber or whether the bridge was always a metaphor for the link between people and gods. Both stories have been told for two thousand years.

Spanish also uses \emph{puente} for a long weekend — the days that bridge a holiday to the nearest Saturday. It is the same idea, doing modern work.
\end{cousinweb}

\begin{grammarlens}[title={What you've built}]
Two, and the way out now crosses something.

\emph{El sendero} and \emph{el puente}. You also have the \emph{ue} break confirmed three times over, which means the next unfamiliar Spanish word with \emph{ue} in the middle is probably a Latin \emph{o} you already know in English.
\end{grammarlens}

\subsection*{Your turn}
Say these aloud:

\begin{itemize}
\raggedright
  \item \emph{el puente}
  \item \emph{el sendero y el puente}
  \item \emph{pueblo}, \emph{cuello}, \emph{puente} — the same break three times
\end{itemize}

\begin{tcolorbox}[breakable,colback=teal!4,colframe=teal!35!black,title={Before you move on}]
What was \emph{pontem}? (A bridge.) What does \emph{pontifex} mean, taken apart? (Bridge-maker.) And which Latin vowel became \emph{ue}? (Short \emph{o}.)
\end{tcolorbox}
\section[la orilla]{la orilla — the sixth little one}

\label{lesson:ES-C295-orilla}



\begin{quote}
Before the new one: what was \emph{pontem}? (A bridge.) And what was \emph{semita}? (A footpath.)
\end{quote}

\subsection*{What to know first}
\begin{itemize}
\raggedright
  \item el lago — one thing that has one.
  \item el puente — what joins two of them.
\end{itemize}

\begin{sounds}
\begin{itemize}
\raggedright
  \item \emph{o-RI-lla}, stress in the middle. The \emph{r} between vowels is a single tap, not rolled. $\to$ reference
\end{itemize}
\end{sounds}

\begin{cousinweb}
\textbf{la orilla} — \textquotedblleft{}the shore,\textquotedblright{} \textquotedblleft{}the bank,\textquotedblright{} \textquotedblleft{}the edge.\textquotedblright{}

Latin \textbf{ora} meant a border or an edge — of a country, of a garment, of the sea. Spanish added the small ending and got \emph{orilla}, \textquotedblleft{}the little edge,\textquotedblright{} and then used it for all three of those things at once.

That breadth is the useful part. The bank of a river is \emph{la orilla}. The shore of a lake is \emph{la orilla}. The hem of a cloth is \emph{la orilla}. One word covers every place where a thing stops being itself, which is a lot of ground for four syllables.

And there is the small ending for the sixth time — an onion, a seed, a knee, an ankle, a knife, and now an edge. Six words with nothing in common except the way they were built. When you meet a new Spanish word ending in \emph{-illa} or \emph{-illo}, take the ending off and look underneath: the plain word is often waiting there, and is often one you already have.
\end{cousinweb}

\begin{grammarlens}[title={What you've built}]
Three, and the route now reaches water: \emph{el sendero}, \emph{el puente}, \emph{la orilla}.

The diminutive has become a genuine tool rather than a piece of trivia. That is worth more than any single one of the six words that taught it to you.
\end{grammarlens}

\subsection*{Your turn}
Say these aloud:

\begin{itemize}
\raggedright
  \item \emph{la orilla}
  \item \emph{la orilla del lago}
  \item \emph{el sendero}, \emph{el puente}, \emph{la orilla}
\end{itemize}

\begin{tcolorbox}[breakable,colback=teal!4,colframe=teal!35!black,title={Before you move on}]
What did \emph{ora} mean? (An edge, a border.) Name three things that have \emph{una orilla}. (A river, a lake, a cloth.) And what does \emph{-illa} add? (Smallness.)
\end{tcolorbox}
\section[la isla]{la isla — and the silent letter English added by mistake}

\label{lesson:ES-C295-isla}



\begin{quote}
Quick pass: what did \emph{ora} mean? (An edge.) What was \emph{pontem}? (A bridge.) And what was \emph{semita}? (A footpath.)
\end{quote}

\subsection*{What to know first}
\begin{itemize}
\raggedright
  \item la orilla — what one is entirely made of.
  \item el agua — what is round it.
\end{itemize}

\begin{sounds}
\begin{itemize}
\raggedright
  \item \emph{IS-la}, two beats, stress on the first. The \emph{s} is clean and unvoiced, never a \emph{z} sound. $\to$ reference
\end{itemize}
\end{sounds}

\begin{cousinweb}
\textbf{la isla} — \textquotedblleft{}the island.\textquotedblright{}

Latin \textbf{insula}, and Spanish simply let the \emph{n} go: \emph{insula}, \emph{isla}. The English words kept it. Something \textbf{insular} has an island's outlook. To \textbf{isolate} is to make an island of somebody — that one came back through Italian \emph{isolato}. A \textbf{peninsula} is \emph{paene insula}, \textquotedblleft{}almost an island.\textquotedblright{} And \textbf{insulin} was named for the little islands of cells in the pancreas that produce it.

English \textbf{isle} has a stranger history. It came from French \emph{ile}, with no \emph{s} in it and none pronounced. Scholars later inserted the \emph{s} to make the word look more like \emph{insula}, and it has sat there silently ever since. Then the same scholars put an \emph{s} into \textbf{island}, which is not from Latin at all and never had one.

So \emph{isla} dropped a letter honestly and \emph{isle} gained one dishonestly, and the two words now look almost the same for opposite reasons.
\end{cousinweb}

\begin{grammarlens}[title={What you've built}]
Four, and the way out has run as far as it can go: \emph{el sendero}, \emph{el puente}, \emph{la orilla}, \emph{la isla}.

Path, crossing, edge, and the thing past the edge. Four words that put you somewhere, and the last of them arrived with four English relatives and a spelling scandal attached.
\end{grammarlens}

\subsection*{Your turn}
Say these aloud:

\begin{itemize}
\raggedright
  \item \emph{la isla}
  \item \emph{la orilla de la isla}
  \item \emph{el sendero}, \emph{el puente}, \emph{la orilla}, \emph{la isla}
\end{itemize}

\begin{tcolorbox}[breakable,colback=teal!4,colframe=teal!35!black,title={Before you move on}]
What was \emph{insula}? (An island.) What does \emph{peninsula} say, taken apart? (Almost an island.) And why does \emph{isle} have an \emph{s}? (Scholars put it there to imitate Latin.)
\end{tcolorbox}
\section[la niebla]{la niebla — the word that was promised}

\label{lesson:ES-C295-niebla}



\begin{quote}
Before the last of the five: what was \emph{insula}? (An island.) What did \emph{ora} mean? (An edge.) And what was \emph{pontem}? (A bridge.)
\end{quote}

\subsection*{What to know first}
\begin{itemize}
\raggedright
  \item la nube — the near-cousin, set aside earlier for this.
  \item la orilla — where it sits thickest.
\end{itemize}

\begin{sounds}
\begin{itemize}
\raggedright
  \item \emph{NIE-bla}, stress on the first part. The \emph{ie} is one slide, and the \emph{b} is soft between the vowel and the \emph{l}. $\to$ reference
\end{itemize}
\end{sounds}

\begin{cousinweb}
\textbf{la niebla} — \textquotedblleft{}the mist,\textquotedblright{} \textquotedblleft{}the fog.\textquotedblright{}

When \emph{la nube} was taken apart, one word was deliberately left on the shelf: \textbf{nebula}, which Latin used for mist as against \emph{nubes} for cloud. This is that word, arrived at last.

The road from \emph{nebula} to \emph{niebla} is short. The short \emph{e} broke into \emph{ie}, as it did in \emph{cielo} and \emph{miel}. The \emph{-bul-} in the middle collapsed to \emph{-bl-}, which Spanish does whenever an unstressed vowel finds itself between two consonants it can be spared from.

English kept \emph{nebula} whole and gave it to astronomy, for the faint clouds between the stars. Something \textbf{nebulous} is misty in the figurative sense — an idea you cannot get an edge on. Which is a good word to end a run on, since the whole point of taking words apart is to make them less nebulous than they were.
\end{cousinweb}

\begin{grammarlens}[title={What you've built}]
Five: \emph{el sendero}, \emph{el puente}, \emph{la orilla}, \emph{la isla}, \emph{la niebla}.

Read in order they are a way out of a place — the path, the crossing, the water's edge, the land beyond it, and the weather closing in behind you.

This run also shuts a circle. It opened with a cloud and closes with the mist that Latin kept carefully separate from it, and the two words have been waiting for each other since. That is the argument for taking a word apart rather than memorising it: the pieces you find keep turning up, and they keep paying.
\end{grammarlens}

\subsection*{Your turn}
Say these aloud:

\begin{itemize}
\raggedright
  \item \emph{la niebla}
  \item \emph{la nube}, then \emph{la niebla} — the cloud and the mist
  \item all five — \emph{el sendero}, \emph{el puente}, \emph{la orilla}, \emph{la isla}, \emph{la niebla}
\end{itemize}

\begin{tcolorbox}[breakable,colback=teal!4,colframe=teal!35!black,title={Before you move on}]
What was \emph{nebula}, as against \emph{nubes}? (Mist, not cloud.) What did the short \emph{e} become? (\emph{Ie}.) And what does \emph{nebulous} mean? (Misty — an idea with no clear edge.)
\end{tcolorbox}
